\documentclass[a4paper,10pt]{article}
\usepackage[francais]{babel}
\usepackage{url}

\input{letter-icube-uds.sty}

\enCouleur

\begin{document}

\name{Stéphane \textsc{Genaud}}
\fonction{Professeur ENSIIE, mise à disposition Université de Strasbourg}
\telephoneillkirch{+33 (0)3688 54542}
\email{genaud@unistra.fr}

\begin{letter}[Recommandation de Garip \textsc{KUSOGLU} pour une offre de thèse.]%
  {à}%
  {Laurent Philippe, Institut FEMTO-ST, UBFC}


\noindent  J'ai été sollicité par Garip KUSOGLU pour cette lettre de recommandation.

  Garip KUSOGLU est  étudiant en deuxième année de  Master Informatique, mention
  \textit{Science de L'Ingénierie  Logicielle} à la faculté  de Mathématiques et
  Informatique  de l'université  de  Strasbourg.  Il  effectue  sa scolarité  en
  Master, depuis la première année, sous un statut d'apprenti. Son apprentissage
  se fait dans le laboratoire ICube sur  un travail de recherche dans mon équipe
  ICPS.  Son maitre  d'apprentissage est  Bérenger Bramas,  chargé de  recherche
  INRIA. Le  rythme d'alternance  dans sa  formation est  en moyenne,  par mois,
  consiste en trois semaines d'apprentissage et 1 semaine de cours académiques.
  
  Je  connais  Garip  depuis  sa  première de  Master  et  j'ai  progressivement
  participé  à son  encadrement sur  son  travail de  recherche. J'ai  également
  cotoyé  récemment Garip  en  tant  qu'enseignant dans  le  cours  de Master  2
  \textit{Traitements Large Echelle} pour les sujet Hadoop et Spark.

À mes yeux, Garip a les qualités nécessaires pour mener à bien un travail de
doctorat. Il  possède des compétences techniques très solides et une très grande
curiosité intellectuelle. Cette curiosité l'a amené à enrichir son bagage de
connnaissances des langages  de programmation et technologies de l'état de l'art
en informatique. De manière plus générale, Garip est curieux des idées
novatrices dans le domaine.

Dans son travail de recherche, il est également très
opiniatre et absolument dédié à son travail. Je peux témoigner qu'une fois
l'objectif fixé, il travaille sans relâche pour atteindre l'objectif. Il a
également les idées très claires sur son travail, et dans nos discussions de
travail, Garip fait preuve d'une grande honnêteté intellectuelle. Ceci nous permet
d'avoir une relation de travail fructueuse.

Etant donné son domaine de compétence lié au parallélisme, et pour les qualités
mentionnées ci-dessus, sa candidature à l'offre \textit{Offline
scheduling in a DZ datacenter with uncertainties} me semble tout à fait pertinente.


  Veuillez agrééer, l'expression de mes sentiments respectueux.

\begin{flushright}
\includegraphics[width=.25\textwidth]{signgenaud.jpg} 
\end{flushright}

\end{letter}


\end{document}
